\setcounter{section}{58} % tema número N-1
\section{Fundamentos de Sistemas Operativos}

Nombre completo: Conceptos y Fundamentos de Sistemas Operativos. Evolución y tendencias.

\localtableofcontents

El sistema Operativo es el componente software mas importante de un ordenador ta que se ocupa de gestionar los recursos de la máquina y ejerce de interfaz entre los componentes hardware y el usuario.

\subsection{Componentes}

\subsubsection{Gestión de procesos}
Los procesos utilizan recursos que proporciona el sistema. También orquestra la sincronización entre procesos. El sistema mantiene una estructura que le permite organizar y describir el estado y propiedad de los recursos, así como ejercer el control sobre las aplicaciones. 

Un SO Multitarea es aquel capaz de ejecutar varios procesos al mismo tiempo. Al contrario que en la multiprogramación, dónde el procesador asigna y desasigna el proceso en marcha.

Los procesos se asignan y desasignas mediante interrupciones software o hardware. Esta orquestación la hace el \textbf{scheduler} según las políticas que implemente, como: Round robin, FIFO, LIFO, SJF, SRT...

Unm proceso puede estar activo, preparado, o bloqueado. El sistema operativo lleva registro de estos estados para cada aplicación. 

\subsubsection{Administración de memoria}

La memoria es usada para almacenar los datos de los programas que se ejecutan en la CPU. El SO es responsable de saber qué partes están libres u ocupadas y cual asignar a quien. También hay varios algoritmos para asignar memoria como: FIFO, NRU, LRU, NFU, LFU, random.

La segmentación es una técnica aplicada a los programas para dividirlos en partes y asignar a cada una un elemento de la memoria. De esta forma, un mismo segmento puede ser compartido por varios programas (como una librería)

\subsubsection{Entrada y Salida}

También se encarga de gestionar la comunicación con los elementos. Pueden ser orientados a bloque (direccionables) o a carácter (no direccionables). 

\begin{description}
    \item[Driver] Parte del sistema encargada de gobernar un dispositivo. Es un software que conoce los detalles de un dispositivo.
    \item[DMA] (Direct Memory Access) Tecnología que permite a un dispositivo transferir datos entre este y la memoria sin pasar por el procesador activamente.
    \item[PnP] Especificación que permite conectar un periférico en caliente sin configurar ningún elemento hardware.
\end{description}

\subsubsection{Gestión de archivos}

Organización que permite almacenar ficheros y datos y usarlos por el sistema y las aplicaciones. Normalmente los dispositivos proveen de una serie de sectores consecutivos, o cadena de bloques. Siendo los sistemas los intermediarios entre esa cadena de bloques y las aplicaciones. También tienen utilidades para gestionar permisos, mover elementos, renombrar, direccionar, agrupar, etc\dots

\begin{itemize}
    \item Sistemas Windows
    \begin{description}
        \item[FAT] desarrollado para DOS. Es sencillo y ampliamente extendido.
        \item[HPFS] Creado para superar las limitaciones de FAT. Permite nombres largos, metadatos, seguridad e información estructural.
        \item[NTFS] Basado en HPFS, permite definir el tamaño de cluster superior a 512B independienteal tamaño de la partición.
    \end{description}
    \item Sistemas UNIX
    \begin{description}
        \item[Ext2] Baja fragmentación, compatible con sistemas grandes de ficheros.
        \item[Ext3] Mejora de Ext2 y compatible. Añade Journaling para posibilitar recuperación frente a fallos.
        \item[Ext4] Mejora de Ext3, menor uso de cpu y mas velocidad de escritura.
        \item[SWAP] Linux utiliza una partición dedicada a la memoria virtual.
    \end{description}
    \item Sistemas OSx
    \begin{description}
        \item[APFS] Estandar moderno, está optimizado para SSD e incluye cifrado.
        \item[HFS+] Estandar antiguo, utilizado aún para discos mecánicos.
    \end{description}

\end{itemize}

\subsection{Tipos de Sistemas Operativos}
Mira macho. Según estas características, una o mucha: usuarios, tareas concurrentes y procesadores.

\subsection{Tendencias}
Cambio en los paradigmas de sistemas con la introducción de los dispositivos móviles.

Google desarrolla Android, basado en Linux. Tiene una cuota de 80\%. Desarrollado en paralelo con Android Open Source Project (Android como tal no es libre)

Apple desarrolla iOS para iPhone, y iPad. Se deriva de macOS, y este de Darwin BSD. Es un sistema fuertemente dependiente del ecosistema de apple.
